\section{Analytical Results}\label{sec:analytical}

This section is organized in two parts.
Section~\ref{subsec:general-nonsep} derives all results under the general
non-separable utility function $U(c^{j},g,n^{j})$ of
Assumption~\ref{asm:general}, working entirely with the elasticities
$(\epsilon^{j}_{cc},\epsilon^{j}_{cg})$. These results hold for any utility
function satisfying Assumption~\ref{asm:general} -- CES, GHH-CRRA, CCRRA, or
any other form -- and require no parametric restrictions.
Section~\ref{subsec:linear-effective} then introduces the linear
effective-consumption form $\tilde c^{j}=c^{j}+\theta^{j}g$
(Assumption~\ref{asm:linear}) as a tractable parameterization that allows
closed-form corollaries, numerical calibration, and simulation. Researchers
wishing to use alternative functional forms -- such as CES utility -- can
substitute their own $(\epsilon^{j}_{cc},\epsilon^{j}_{cg})$ directly into the
general results of Section~\ref{subsec:general-nonsep}.

Full derivations are in Appendix~\ref{app:general-nonsep}; the complete model
setup is in Appendix~\ref{app:setup}.

\subsection{General non-separable preferences}\label{subsec:general-nonsep}

\subsubsection{Two sufficient statistics}
The standard separable HANK multiplier depends on one sufficient statistic:
income cyclicality $\chi$ \citep{Bilbiie2025}. Introducing $U_{cg}^{j}\neq 0$
creates a second: the direct fiscal channel $\xi$. These two objects are
cleanly separated. $\chi$ is determined by own-consumption curvature
$\epsilon^{j}_{cc}$ only. $\xi$ is determined by the cross-elasticities
$\epsilon^{j}_{cg}$ only. This separation holds in full generality under
Assumption~\ref{asm:general} and is the structural foundation of all results
that follow.

\subsubsection{Notation: tax progressivity}
Following \citet{Bilbiie2025}, we implement a balanced budget $T_{t}=G_{t}$ with
profit redistribution rate $\tau_{D}$ and tax progressivity parameter $\mu$ (we
use $\mu$ rather than \bilbiie{}'s $\alpha$ to avoid clash with
$\alpha^{j}\equiv\bar c/\bar{\tilde c}^{j}$). \HtM{} households pay per-capita
taxes $\mu/\lambda$ per unit of spending. Since $T_{t}=G_{t}$, $\hat g_{t}$
appears in the $H$ budget constraint as a tax term: the balanced budget forces
taxes to equal spending, and $H$ households bear fraction $\mu/\lambda$.

\subsubsection{Step 1: intratemporal conditions}
From \eqref{eq:intratemp}:
\begin{equation}\label{eq:intratemp-restate}
  \phi\,\hat n^{j}_{t}=\hat w_{t}-\epsilon^{j}_{cc}\hat c^{j}_{t}
  -\epsilon^{j}_{cg}\hat g_{t},\qquad j\in\{S,H\}.
\end{equation}
Under Assumption~\ref{asm:linear}: $\epsilon^{j}_{cc}=\sigma\alpha^{j}$,
$\epsilon^{j}_{cg}=\sigma\Gamma^{j}$.

\subsubsection{Step 2: the real wage}
Aggregating \eqref{eq:intratemp-restate} and using
$\hat c^{S}_{t}=(\hat c_{t}-\lambda\hat c^{H}_{t})/(1-\lambda)$:
\begin{equation}\label{eq:wage-1}
  \hat w_{t}=(\phi+\epsilon^{S}_{cc})\hat c_{t}
    +\lambda(\epsilon^{H}_{cc}-\epsilon^{S}_{cc})\hat c^{H}_{t}
    +\bigl[\phi+\lambda\epsilon^{H}_{cg}+(1-\lambda)\epsilon^{S}_{cg}\bigr]\hat g_{t}.
\end{equation}
Writing $\hat w_{t}=W_{c}\hat c_{t}+W_{g}\hat g_{t}$ with
$\hat c^{H}_{t}=\chi\hat c_{t}+\xi\hat g_{t}$:
\begin{align}
  W_{c}&\equiv(\phi+\epsilon^{S}_{cc})+\lambda\chi(\epsilon^{H}_{cc}-\epsilon^{S}_{cc}),
    \label{eq:Wc}\\
  W_{g}&\equiv\phi+\lambda\epsilon^{H}_{cg}+(1-\lambda)\epsilon^{S}_{cg}
    +\lambda\xi(\epsilon^{H}_{cc}-\epsilon^{S}_{cc}).\label{eq:Wg}
\end{align}

\subsubsection{Step 3: income cyclicality \texorpdfstring{$\chi$}{chi}}
Define $K\equiv\phi+1-\phi\tau_{D}/\lambda$ and
$D\equiv(\phi+\epsilon^{H}_{cc})-K\lambda(\epsilon^{H}_{cc}-\epsilon^{S}_{cc})$.

\begin{proposition}[Income cyclicality -- general]\label{prop:chi-general}
Under Assumption~\ref{asm:general},
\begin{equation}\label{eq:chi-general}
  \chi=\frac{K(\phi+\epsilon^{S}_{cc})}{D}.
\end{equation}
$\chi$ depends on $\epsilon^{S}_{cc}$ and $\epsilon^{H}_{cc}$ but not on
$\epsilon^{j}_{cg}$: income cyclicality is determined by own-consumption
curvature, not by the non-separability between $c$ and $g$.
Under Assumption~\ref{asm:linear}:
$D=(\phi+\sigma\alpha^{H})-K\sigma\lambda(\alpha^{H}-\alpha^{S})$ and
\begin{equation}\label{eq:chi-linear}
  \chi(\theta^{S},\theta^{H})=\frac{K(\phi+\sigma\alpha^{S})}{D}.
\end{equation}
When $\theta^{S}=\theta^{H}$: $\chi=K=1+\phi(1-\tau_{D}/\lambda)$, \bilbiie{}'s
expression, independent of $\theta$.
\end{proposition}

\begin{remark}[Why $\epsilon^{j}_{cg}$ does not appear in $\chi$]
Income cyclicality $\chi$ measures how \HtM{} income co-moves with aggregate
output. This is determined by how wage income responds to aggregate demand --
governed by $\epsilon^{j}_{cc}$ through the labor-supply elasticity -- not by
how directly government spending affects marginal utility. The non-separability
between $c$ and $g$ enters the direct fiscal channel $\xi$, not $\chi$.
\end{remark}

\subsubsection{Step 4: direct fiscal channel \texorpdfstring{$\xi$}{xi}}

\begin{proposition}[Direct fiscal channel -- general]\label{prop:xi-general}
Under Assumption~\ref{asm:general},
\begin{equation}\label{eq:xi-general}
  \xi=\frac{K\bigl[\phi+\lambda\epsilon^{H}_{cg}+(1-\lambda)\epsilon^{S}_{cg}\bigr]
          -\epsilon^{H}_{cg}-\mu\phi/\lambda}{D}.
\end{equation}
$\xi$ depends on both $\epsilon^{H}_{cg}$ and $\epsilon^{S}_{cg}$. Under
separability ($\epsilon^{j}_{cg}=0$): $\xi\rightarrow\zeta(\chi-\mu/\lambda)$,
\bilbiie{}'s expression.
Under Assumption~\ref{asm:linear}, $\epsilon^{j}_{cg}=\sigma\Gamma^{j}$:
\begin{equation}\label{eq:xi-linear}
  \xi(\theta^{S},\theta^{H})=
  \frac{K\bigl[\phi+\sigma\lambda\Gamma^{H}+\sigma(1-\lambda)\Gamma^{S}\bigr]
        -\sigma\Gamma^{H}-\mu\phi/\lambda}{D}.
\end{equation}
\end{proposition}

\noindent In the separable case $\theta^{S}=\theta^{H}=0$ (so $\Gamma^{j}=0$,
$\alpha^{j}=1$, $D=\phi+\sigma$, $K=\chi$):
$\xi(0,0)=\zeta(\chi-\mu/\lambda)$ where $\zeta\equiv\phi/(\phi+\sigma)$,
matching \citet{Bilbiie2025} exactly.

\subsubsection{Step 5: the saver Euler equation}
The saver's Euler equation with type-specific $\theta^{S}$:
\begin{equation}\label{eq:euler-linear}
  \hat c^{S}_{t}=\mathbb{E}_{t}\hat c^{S}_{t+1}
    -\frac{1}{\sigma\alpha^{S}}\hat r_{t}
    -\frac{\Gamma^{S}}{\alpha^{S}}(\hat g_{t}-\mathbb{E}_{t}\hat g_{t+1}).
\end{equation}
When $\theta^{S}<0$ (complements): $\Gamma^{S}<0$, so the last term is positive
when $\hat g_{t}>\mathbb{E}_{t}\hat g_{t+1}$ -- a direct demand stimulus. When
$\theta^{S}>0$ (substitutes): the term is negative -- government spending
crowds out saver consumption. In THANK, the self-insurance Euler equation
introduces the compounding coefficient
\begin{equation}\label{eq:delta}
  \delta\equiv 1+(\chi-1)\frac{1-s}{1-\lambda\chi},
\end{equation}
which depends on $\chi$ but not directly on $\theta^{j}$.

\subsubsection{Step 6: the planned expenditure curve and multiplier}
The PE curve in THANK with general $(\theta^{S},\theta^{H})$:
\begin{equation}\label{eq:PE}
  \hat c_{t}=[1-\beta(1-\lambda\chi)]\hat y_{t}
    -(1-\lambda)\frac{\beta\sigma}{\alpha^{S}}\hat r_{t}
    +\beta\delta(1-\lambda\chi)\mathbb{E}_{t}\hat c_{t+1}
    +\beta\lambda\xi(\theta^{S},\theta^{H})(\hat g_{t}-\mathbb{E}_{t}\hat g_{t+1}).
\end{equation}

Setting $\hat r_{t}=0$:

\begin{proposition}[General fiscal multiplier]\label{prop:multiplier-general}
Under Assumption~\ref{asm:general}, the on-impact fiscal multiplier in THANK
with heterogeneous non-separable preferences is
\begin{equation}\label{eq:multiplier}
  \mu=1+\frac{\lambda\xi}{1-\lambda\chi},
\end{equation}
where $\chi$ and $\xi$ are given by \eqref{eq:chi-general} and
\eqref{eq:xi-general}, and
$D=(\phi+\epsilon^{H}_{cc})-K\lambda(\epsilon^{H}_{cc}-\epsilon^{S}_{cc})$.
\end{proposition}

\begin{proof}
With $\hat r_{t}=0$, $\hat y_{t}=\hat c_{t}+\hat g_{t}$, and
$\hat c^{H}_{t}=\chi\hat c_{t}+\xi\hat g_{t}$, substitute into goods-market
clearing. Collecting terms: $(1-\lambda\chi)\hat y_{t}=(1-\lambda\chi+\lambda\xi)\hat g_{t}$.
Dividing by $(1-\lambda\chi)$ yields \eqref{eq:multiplier}.
\end{proof}

\begin{corollary}[General heterogeneous internalization]\label{cor:het}
Under Assumption~\ref{asm:general} with
$\epsilon^{H}_{cg}\neq\epsilon^{S}_{cg}$, the direct fiscal channel is given by
\eqref{eq:xi-general}. The multiplier exceeds the separable benchmark
($\xi>\xi^{\mathrm{sep}}$) if and only if the \emph{distributional dominance}
condition holds:
\begin{equation}\label{eq:dom-cond}
  (K\lambda-1)\epsilon^{H}_{cg}+K(1-\lambda)\epsilon^{S}_{cg}>0.
\end{equation}
When $\epsilon^{H}_{cg}<0$ (\HtM{} complements, $U^{H}_{cg}>0$) and
$\epsilon^{S}_{cg}>0$ (savers substitute, $U^{S}_{cg}<0$), the two terms push
in opposite directions. Condition \eqref{eq:dom-cond} is satisfied when \HtM{}
complementarity is sufficiently strong relative to saver substitutability,
weighted by $K\lambda-1$.
\end{corollary}

\begin{proposition}[General MPC/MPE decoupling]\label{prop:decoupling}
Under Assumption~\ref{asm:general} with $\epsilon^{H}_{cc}\neq\epsilon^{S}_{cc}$,
income cyclicality $\chi$ depends on $\epsilon^{H}_{cc}-\epsilon^{S}_{cc}$
independently of $\lambda$, while the aggregate MPC depends on $\lambda$ but
not on $\epsilon^{j}_{cc}$. A researcher can therefore independently target
\textbf{(i)} the aggregate MPC, by choosing $\lambda$;
\textbf{(ii)} the aggregate MPE ($=\chi$), by choosing
$\epsilon^{H}_{cc}-\epsilon^{S}_{cc}$;
\textbf{(iii)} the fiscal multiplier, by choosing
$\epsilon^{H}_{cg}-\epsilon^{S}_{cg}$ through \eqref{eq:dom-cond}.
In the general specification these three instruments are mutually independent:
$\epsilon^{j}_{cc}$ and $\epsilon^{j}_{cg}$ are free parameters not linked by
any parametric restriction.
\end{proposition}

\begin{remark}[Applicability to alternative functional forms]
Propositions~\ref{prop:multiplier-general} and~\ref{prop:decoupling} and
Corollary~\ref{cor:het} apply to any utility function satisfying
Assumption~\ref{asm:general}. For example, a researcher using CES utility
$U^{j}=\bigl[(c^{j})^{\rho}+\psi(g)^{\rho}\bigr]^{(1-\sigma)/\rho}$ need only
compute the steady-state elasticities $\epsilon^{j}_{cc}$ and $\epsilon^{j}_{cg}$
from their chosen functional form and substitute directly into
\eqref{eq:chi-general}, \eqref{eq:xi-general}, and \eqref{eq:multiplier}. No
re-derivation is required.
\end{remark}

\subsection{Linear effective consumption: a tractable
parameterization}\label{subsec:linear-effective}
We now specialize the general results to the linear effective consumption form
of Assumption~\ref{asm:linear}: $\tilde c^{j}=c^{j}+\theta^{j}g$, which gives
$\epsilon^{j}_{cc}=\sigma\alpha^{j}$ and $\epsilon^{j}_{cg}=\sigma\Gamma^{j}$
where $\alpha^{j}=\bar c/(\bar c+\theta^{j}\bar g)$ and
$\Gamma^{j}=\theta^{j}\bar g/(\bar c+\theta^{j}\bar g)$. This parameterization
is convenient for three reasons: (i) it collapses heterogeneous internalization
to a single parameter $\theta^{j}$ per type; (ii) it delivers closed-form
corollaries; and (iii) it maps directly to calibration and simulation. The
parameter $\theta^{j}\in(-\bar c/\bar g,\infty)$ with $\theta^{j}=0$ nesting
separability, $\theta^{j}<0$ complementarity, and $\theta^{j}>0$
substitutability.

\begin{corollary}[\bilbiie{} separable benchmark]\label{cor:bilbiie}
When $\theta^{S}=\theta^{H}=0$: $\chi=1+\phi(1-\tau_{D}/\lambda)$,
$\xi=\zeta(\chi-\mu/\lambda)$, and
\begin{equation}\label{eq:mult-bilbiie}
  \mu^{\mathrm{Bilbiie}}=1+\frac{\lambda\zeta(\chi-\mu/\lambda)}{1-\lambda\chi}.
\end{equation}
\end{corollary}

\begin{corollary}[RANK benchmark]\label{cor:rank}
In the representative-agent limit $\lambda\to 0$, for any $(\theta^{S},\theta^{H})$:
\begin{equation}
  \mu^{\mathrm{RANK}}(\theta^{S})=1-\psi^{S}=1+\frac{|\theta^{S}|\bar g}{\bar c}.
\end{equation}
Only the saver direct-demand channel survives. This recovers
\citet{BouakezRebei2007} under $\theta^{S}<0$.
\end{corollary}

\begin{corollary}[Symmetric complementarity]\label{cor:sym}
When $\theta^{S}=\theta^{H}=\theta<0$: $\alpha^{S}=\alpha^{H}=\alpha>1$,
$D=\phi+\sigma\alpha$, $\chi=K=1+\phi(1-\tau_{D}/\lambda)$ (identical to
\bilbiie{}'s separable expression), and
\begin{equation}\label{eq:xi-sym}
  \xi(\theta)=\zeta_{\alpha}\Bigl(\chi-\tfrac{\mu}{\lambda}+\tfrac{\sigma\Gamma(\chi-1)}{\phi}\Bigr),
\end{equation}
where $\zeta_{\alpha}=\phi/(\phi+\sigma\alpha)<\zeta=\phi/(\phi+\sigma)$.
For $\chi>1$ (countercyclical inequality, empirically standard):
$\xi(\theta)<\xi(0)$ because (i) $\zeta_{\alpha}<\zeta$ (scaling dampening)
and (ii) $\sigma\Gamma(\chi-1)/\phi<0$ (since $\Gamma<0$ under complementarity).
Consequently,
$\mu^{\mathrm{sym}}(\theta)<\mu^{\mathrm{Bilbiie}}$ when $\chi>1$.
\end{corollary}

\begin{corollary}[Heterogeneous internalization]\label{cor:het-explicit}
Under the empirically plausible case $\theta^{S}>0$ (savers: substitutes) and
$\theta^{H}<0$ (\HtM: complements): $\alpha^{S}<1<\alpha^{H}$,
$\Gamma^{S}>0>\Gamma^{H}$. Income cyclicality is
\begin{equation}
  \chi(\theta^{S},\theta^{H})=
  \frac{K(\phi+\sigma\alpha^{S})}
       {(\phi+\sigma\alpha^{H})-K\sigma\lambda(\alpha^{H}-\alpha^{S})}.
\end{equation}
Since $\alpha^{H}>\alpha^{S}$, the denominator $D<\phi+\sigma\alpha^{H}$,
making $\chi$ larger than in the symmetric case -- heterogeneous internalization
\emph{raises} income cyclicality. The multiplier exceeds the separable
benchmark iff
\begin{equation}\label{eq:dom-threshold}
  \underbrace{\frac{\lambda|\xi(\theta^{S},\theta^{H})|}{1-\lambda\chi}}_{\text{HtM amplification}}
  >\underbrace{|\psi^{S}|\cdot(1-\lambda)}_{\text{Saver dampening}}.
\end{equation}
The left side is increasing in $\lambda$; the right side is decreasing. There
exists a threshold $\lambda^{*}$ above which complementarity dominates.
\end{corollary}

\begin{remark}[Fiscal policy and inequality]\label{rem:ineq-amp}
Corollary~\ref{cor:het-explicit} implies that fiscal policy is more effective
in more unequal economies through two reinforcing channels. The standard HANK
channel: more \HtM{} households means higher aggregate MPC and stronger NK
Cross. Our new channel: \HtM{} households are precisely those for whom
government spending is most complementary ($\theta^{H}<0$), raising $\chi$
and $\xi$ beyond the symmetric case. Inequality and complementarity are
co-amplifying forces -- not just additive.
\end{remark}

\subsection{Relation to the MPC/MPE trilemma}\label{subsec:trilemma}

The MPC/MPE trilemma was identified by \citet{Auclert2024}: New Keynesian
models with frictionless labor supply cannot simultaneously match high
empirical MPCs, low MPEs, and moderate fiscal multipliers. Their resolution
requires sticky wages. \BHL{} provide an alternative resolution through $c$--$n$
complementarity ($U_{cn}\neq 0$), which separates income effects on labor
supply from the intertemporal substitution elasticity. This subsection shows
that our $c$--$g$ non-separability ($U_{cg}^{j}\neq 0$ with
$\theta^{H}\neq\theta^{S}$) provides a third, orthogonal resolution -- one that
operates through the government-spending margin and requires no modification
to the $c$--$n$ structure.

\begin{proposition}[MPC/MPE decoupling under heterogeneous internalization]
\label{prop:mpc-mpe-decoupling}
Under Assumption~\ref{asm:linear} with $\theta^{S}\neq\theta^{H}$, the
aggregate MPC (governed primarily by $\lambda$) and the aggregate MPE (governed
by $\chi(\theta^{S},\theta^{H})$) can be independently targeted. The mechanism
is $\theta$-induced curvature heterogeneity: since $\epsilon^{j}_{cc}=\sigma\alpha^{j}$
under Assumption~\ref{asm:linear}, heterogeneous $\theta^{j}$ induces
$\alpha^{H}\neq\alpha^{S}$, which shifts $\chi$ through the denominator $D$
independently of $\lambda$. Specifically:
\textbf{(i)} $\partial\chi/\partial\lambda\neq 0$ and
$\partial(\mathrm{MPC})/\partial\lambda\neq 0$ -- both are affected by $\lambda$.
\textbf{(ii)} $\partial\chi/\partial\alpha^{S}\neq 0$ generically, while
$\partial(\mathrm{MPC})/\partial\alpha^{S}=0$: the curvature gap
$\alpha^{H}-\alpha^{S}$ (induced by $\theta^{H}-\theta^{S}$) shifts income
cyclicality (MPE) without affecting the aggregate MPC.
Under symmetry ($\theta^{S}=\theta^{H}$): $\alpha^{H}=\alpha^{S}$,
$D=\phi+\sigma\alpha$, $\chi=K$ -- identical to \bilbiie{}'s separable
expression. The decoupling disappears and the trilemma is not resolved within
this case.
\end{proposition}

\begin{table}[ht]
\centering
\small
\caption{Three complementary resolutions of the MPC/MPE trilemma.}
\label{tab:trilemma-comparison}
\begin{threeparttable}
\begin{tabular}{@{}p{2.8cm}p{3.0cm}p{3.4cm}p{3.4cm}@{}}
\toprule
 & \citet{Auclert2024} & \BHL{} & \textbf{This paper} \\
\midrule
Utility & $U(c,n)$ & $U(c,n)$, $U_{cn}\neq 0$ & $U(c,g,n)$, $U_{cg}^{j}\neq 0$ \\[3pt]
Non-sep.\ type & None (baseline) & $c$--$n$ (cons--lab) & $c$--$g$ (cons--govt) \\[3pt]
Trilemma route & Sticky wages & $c$--$n$ curvature & $\theta^{H}-\theta^{S}$ gap \\[3pt]
Requires het.?  & No & No & Yes ($\theta^{H}\neq\theta^{S}$) \\[3pt]
Fiscal question & Multipliers & Multipliers + FG & Multiplier only \\[3pt]
Forward guidance & Not addressed & Catch-22 resolved & Outside scope \\
\bottomrule
\end{tabular}
\begin{tablenotes}\footnotesize
\item The three mechanisms are mutually compatible: a model combining sticky
wages, $U_{cn}\neq 0$, and $U_{cg}^{j}\neq 0$ would have the maximum flexibility
to match all three facts simultaneously. We isolate the $c$--$g$ channel in its
purest form by maintaining flexible wages and separable $c$--$n$ preferences.
\end{tablenotes}
\end{threeparttable}
\end{table}

\subsection{The HANK puzzles: where our framework stands}\label{subsec:puzzles}

The HANK literature has identified several key challenges -- tensions between
empirical targets and model implications -- that any tractable HANK framework
must confront. We now systematically address each.

\paragraph{Puzzle 1: the MPC/MPE trilemma.}
\emph{Resolution}: Under $\theta^{H}\neq\theta^{S}$, the gap
$\theta^{H}-\theta^{S}$ shifts $\chi$ independently of $\lambda$
(Proposition~\ref{prop:mpc-mpe-decoupling}). This works even with fully
separable $c$--$n$ preferences: the resolution operates entirely through the
government-spending margin.

\paragraph{Puzzle 2: the forward guidance puzzle / Catch-22.}
\emph{Outside scope}. Our framework is complementary to the Catch-22 resolution:
once the forward-guidance puzzle is addressed (e.g.\ by \bilbiie{}'s risk
channel), our $c$--$g$ complementarity layer can be added on top to further
discipline fiscal multipliers.

\paragraph{Puzzle 3: the crowding-out puzzle.}
\emph{Directly resolved and generalized to HANK}. Under complementarity
($\epsilon^{j}_{cg}<0$), the fiscal multiplier exceeds the separable benchmark
for any $\lambda>0$ through both the saver direct-demand channel (Euler
equation) and the \HtM{} labor-supply channel. Our framework is the first to
show this crowding-out resolution in a tractable heterogeneous-agent model.

\subsection{TANK vs.\ THANK}\label{subsec:tank-thank}
TANK is the special case $s=1$ (no switching), giving $\delta=1$. The on-impact
multiplier \eqref{eq:multiplier} is identical in TANK and THANK because
$\delta$ multiplies only $\mathbb{E}_{t}\hat c_{t+1}$, which is irrelevant for
transitory shocks. For persistent shocks ($\rho_{g}>0$), THANK generates
additional amplification when $\chi>1$ (countercyclical inequality) through
$\delta>1$ -- Euler-equation compounding from precautionary saving. Note
$\delta$ does not depend on $\theta^{j}$ since $\chi$ in the general case
depends on $(\theta^{S},\theta^{H})$ but $\delta$ is a function of $\chi$, $s$,
and $\lambda$ only.
